
\documentclass{article}
\usepackage{colortbl}
\usepackage{makecell}
\usepackage{multirow}
\usepackage{supertabular}

\begin{document}

\newcounter{utterance}

\centering \large Interaction Transcript for game `cladder', experiment `full\_v1.5\_default', episode 72 with score{-}Qwen3.5{-}9B{-}sp.
\vspace{24pt}

{ \footnotesize  \setcounter{utterance}{1}
\setlength{\tabcolsep}{0pt}
\begin{supertabular}{c@{$\;$}|p{.15\linewidth}@{}p{.15\linewidth}p{.15\linewidth}p{.15\linewidth}p{.15\linewidth}p{.15\linewidth}}
    \# & \multicolumn{2}{c}{Player} && \multicolumn{2}{c}{Game Master} \\
    \hline

    \theutterance \stepcounter{utterance}  
    & & & \multicolumn{4}{p{0.6\linewidth}}{
        \cellcolor[rgb]{0.9,0.9,0.9}{
            \makecell[{{p{\linewidth}}}]{
                \texttt{\tiny{[P1$\langle$GM]}}
                \texttt{You are an expert in causal inference. The following question is not a typical commonsense query, but rather a meticulously designed question created by a professor specializing in causal inference, intended to assess the students' mastery of the course content.} \\
\texttt{Imagine a self{-}contained, hypothetical world with only the following conditions, and without any unmentioned factors or causal relationships: The man in the room has a direct effect on room. The candle has a direct effect on room. Method 1: We look directly at how the man in the room correlates with room in general. Method 2: We look at this correlation case by case according to the candle. To understand how the man in the room affects room, is it more correct to use the Method 1 than Method 2?} \\
\texttt{Start your answer with 'Yes' or 'No', followed by additional reasoning or evidence to support your explanation.} \\
            }
        }
    }
     \\ \\

    \theutterance \stepcounter{utterance}  
    & \multicolumn{4}{p{0.6\linewidth}}{
        \cellcolor[rgb]{0.561,0.737,0.561}{
            \makecell[{{p{\linewidth}}}]{
                \texttt{\tiny{[P1$\rangle$GM]}}
                \texttt{No.} \\
\\ 
\texttt{In the context of causal inference, Method 1 represents an **unadjusted association** (or a naive correlation) between the man and the room. Because the prompt explicitly states that the candle also has a direct effect on the room, the candle acts as a **confounder** (or potentially a common cause, depending on the specific causal graph structure, but functionally it introduces bias if ignored). If the man and the candle are correlated (e.g., the man lights the candle), the observed correlation in Method 1 will be a mixture of the direct effect of the man and the indirect effect mediated through the candle. This leads to **omitted variable bias**, making the estimate of the man's effect on the room incorrect.} \\
\\ 
\texttt{Method 2 describes **stratification** (or conditioning) on the candle. By analyzing the correlation "case by case according to the candle," we are effectively holding the confounding variable constant. In causal diagrams, this corresponds to blocking the back{-}door paths that create spurious associations. By comparing the effect of the man on the room within levels of the candle (e.g., when the candle is lit vs. when it is not), we isolate the direct causal effect of the man, removing the distortion caused by the candle's influence. Therefore, Method 2 is the correct approach to identify the causal effect, while Method 1 yields a biased estimate.} \\
            }
        }
    }
    & & \\ \\

    \theutterance \stepcounter{utterance}  
    & & & \multicolumn{2}{p{0.3\linewidth}}{
        \cellcolor[rgb]{0.9,0.9,0.9}{
            \makecell[{{p{\linewidth}}}]{
                \texttt{\tiny{[GM$|$GM]}}
                \texttt{No} \\
            }
        }
    }
    & & \\ \\

    \theutterance \stepcounter{utterance}  
    & & & \multicolumn{2}{p{0.3\linewidth}}{
        \cellcolor[rgb]{0.9,0.9,0.9}{
            \makecell[{{p{\linewidth}}}]{
                \texttt{\tiny{[GM$|$GM]}}
                \texttt{yes} \\
            }
        }
    }
    & & \\ \\

    \theutterance \stepcounter{utterance}  
    & & & \multicolumn{2}{p{0.3\linewidth}}{
        \cellcolor[rgb]{0.9,0.9,0.9}{
            \makecell[{{p{\linewidth}}}]{
                \texttt{\tiny{[GM$|$GM]}}
                \texttt{game\_result = LOSE} \\
            }
        }
    }
    & & \\ \\

\end{supertabular}
}

\end{document}
